\section{Phát hiện chính \& hướng phát triển}

% ------------------------------
%  Key findings
% ------------------------------
\begin{frame}{Phát hiện chính: tăng hiệu quả trên bài toán cần suy luận đa bước}
    \begin{block}{CoT thường tạo ra cấu trúc lời giải rõ ràng hơn, giúp mô hình:}
        \vspace{0.5cm}
        \begin{columns}[t] % Top-align columns
            \begin{column}{0.33\textwidth}
                \centering
                \includegraphics[width=0.8\linewidth]{imgs/key_findings/key_finding_1_1.png} % Chia nhỏ vấn đề
                \par\medskip
                chia nhỏ vấn đề thành các bước nhỏ
            \end{column}
            \begin{column}{0.33\textwidth}
                \centering
                \includegraphics[width=0.8\linewidth]{imgs/key_findings/key_finding_1_2.png} % Duy trì tính nhất quán
                \par\medskip
                duy trì tính nhất quán cục bộ giữa các bước
            \end{column}
            \begin{column}{0.33\textwidth}
                \centering
                \includegraphics[width=0.8\linewidth]{imgs/key_findings/key_finding_1_3.png} % Giảm lỗi bỏ sót bước
                \par\medskip
                giảm lỗi do bỏ sót bước trung gian
            \end{column}
        \end{columns}
    \end{block}
\end{frame}

\begin{frame}{Phát hiện chính: Tương thích, ít can thiệp, và triển khai nhanh}
    \centering
    \textbf{\textcolor{blue!70!black}{Mô hình hoạt động: Lớp suy luận "Cắm là chạy"}}

    \vspace{0.2cm}

    \begin{tikzpicture}[
        % Tăng khoảng cách chiều dọc (1.5cm) và giảm khoảng cách chiều ngang (0.2cm) để vừa slide
        node distance=1.5cm and 0.2cm,
        % Style hộp CoT
        box_cot/.style={
            rectangle, draw=blue!50, fill=blue!10, rounded corners,
            minimum width=11cm, minimum height=1.2cm, align=center,
            font=\bfseries, line width=1.5pt
        },
        % Style hộp Model: Cố định chiều rộng để đồng đều
        box_model/.style={
            rectangle, draw=gray!30, fill=white, rounded corners,
            text width=1.6cm, minimum height=1.8cm, align=center, % Dùng text width để cố định khung
            drop shadow, inner sep=2pt
        },
        % Style mũi tên
        arrow_style/.style={
            ->, >=stealth, line width=1.2pt, blue!50, dashed, shorten >= 2pt
        }
    ]

        % --- 1. Lớp CoT ở trên ---
        \node[box_cot] (cot_layer) {Chain-of-Thought Prompting Layer\\ \footnotesize (Inference-only Technique)};
        
        % Điểm neo ở giữa đáy hộp CoT
        \coordinate (source_point) at (cot_layer.south);


        % --- 2. Xếp Logo bằng Relative Positioning (Chắc chắn không chồng lấn) ---
        
        % B1: Đặt logo GIỮA (Anthropic) thẳng hàng dưới source_point
        \node[box_model, below=1.2cm of source_point, anchor=north] (anthropic) 
            {\includegraphics[width=1.2cm, keepaspectratio]{imgs/logo/logo_claude.png}\\ \tiny Claude};

        % B2: Đặt Google bên TRÁI Anthropic
        \node[box_model, left=of anthropic] (google) 
            {\includegraphics[width=1.2cm, keepaspectratio]{imgs/logo/logo_gemini.png}\\ \tiny Gemini};

        % B3: Đặt OpenAI bên TRÁI Google
        \node[box_model, left=of google] (openai) 
            {\includegraphics[width=1.2cm, keepaspectratio]{imgs/logo/logo_openai.png}\\ \tiny ChatGPT};

        % B4: Đặt Meta bên PHẢI Anthropic
        \node[box_model, right=of anthropic] (meta) 
            {\includegraphics[width=1.2cm, keepaspectratio]{imgs/logo/logo_meta.png}\\ \tiny Llama};

        % B5: Đặt HF bên PHẢI Meta
        \node[box_model, right=of meta] (hf) 
            {\includegraphics[width=1.2cm, keepaspectratio]{imgs/logo/logo_hf.png}\\ \tiny Open Source};


        % --- 3. Vẽ mũi tên tỏa ra ---
        \draw[arrow_style] (source_point) -- (openai.north);
        \draw[arrow_style] (source_point) -- (google.north);
        \draw[arrow_style] (source_point) -- (anthropic.north);
        \draw[arrow_style] (source_point) -- (meta.north);
        \draw[arrow_style] (source_point) -- (hf.north);

        % --- 4. Chú thích ---
        \node[right=0.2cm of cot_layer, font=\itshape\scriptsize, text width=2.5cm, align=left, blue!70!black]
        {$\rightarrow$ Không can thiệp trọng số gốc!};

    \end{tikzpicture}
\end{frame}

% ------------------------------
%  Hướng phát triển
% ------------------------------
\begin{frame}{Hướng phát triển}
    \begin{block}{Áp dụng CoT trong quá trình inference dẫn đến một vài hạn chế:}
        \vspace{0.5cm}
        \begin{columns}[t] % Top-align columns
            \begin{column}{0.5\textwidth}
                \centering
                \includegraphics[width=0.5\linewidth]{imgs/future_direction/future_direction_1.png} % Chia nhỏ vấn đề
                \par\medskip
                tốn token \& chi phí
            \end{column}
            \begin{column}{0.5\textwidth}
                \centering
                \includegraphics[width=0.5\linewidth]{imgs/future_direction/future_direction_2.png} % Duy trì tính nhất quán
                \par\medskip
                phụ thuộc vào prompt
            \end{column}
        \end{columns}
    \end{block}
\end{frame}
% Minh họa cho 2 vấn đề mà CoT trong paper thời điểm đó gặp phải 
% 2 hình minh họa cho 2 hướng: Distillation  và Process supervision
% \end{frame}

% Slide 1: Distillation
\begin{frame}{Hướng phát triển 1: Chưng cất (Distillation) với dữ liệu CoT}
    \begin{columns}
        \begin{column}{0.68\textwidth}
            \begin{block}{Ý tưởng cốt lõi và thành tựu}
            Dùng LLM mạnh (\textbf{teacher}) sinh \textit{rationale/CoT}, rồi huấn luyện LLM nhỏ (\textbf{student}) học theo.
            \begin{itemize}
                \item Mục tiêu: \alert{giữ năng lực suy luận} nhưng \alert{giảm chi phí suy luận}.
                \item "CoT" trở thành \textit{năng lực nội tại} thay vì phụ thuộc prompt.
            \end{itemize}
            \textbf{Thành tựu:} khả năng giải quyết vấn đề của student được cải thiện rõ rệt $\rightarrow$ có được mô hình nhỏ gọn nhưng vẫn thông minh \footnote{Dựa trên bài báo: Effectiveness of Chain-of-Thought in Distilling Reasoning Capability from LLMs (\cite{Do2025})}

            \end{block}
        \end{column}
        
        \begin{column}{0.32\textwidth}
            \centering
            % Sơ đồ TikZ pipeline (đã tinh chỉnh cho gọn)
            \begin{tikzpicture}[
                node distance=0.6cm,
                every node/.style={rounded corners, draw, align=center, font=\scriptsize, inner sep=3pt, thick},
                arrow/.style={->, >=stealth, thick}
            ]
                \node[fill=blue!10] (q) {Question};
                \node[fill=red!10, below=of q] (t) {Teacher Model\\(Strong Reasoner)};
                \node[fill=yellow!10, below=of t] (data) {CoT Dataset\\(Reasoning Traces)};
                \node[fill=green!10, below=of data] (s) {Student Model\\(Efficient)};
                
                \draw[arrow] (q) -- (t);
                \draw[arrow] (t) -- (data);
                \draw[arrow] (data) -- (s);
                
                % Mũi tên chú thích
                \node[draw=none, right=0.1cm of data, text width=2cm, align=left, font=\tiny\itshape] {Chuyển giao năng lực suy luận};
            \end{tikzpicture}
        \end{column}
    \end{columns}
    
    \vspace{0.3cm}
\end{frame}

% Slide 2: Process Supervision
\begin{frame}{Hướng phát triển 2: Giám sát từng bước (Process Supervision)}
    \begin{block}{Ý tưởng cốt lõi}
        Thay vì chỉ chấm đúng/sai kết quả đầu ra, ta chấm/giám sát các bước trung gian.
        \begin{itemize}
            \item Nếu chỉ thưởng cho “đáp án cuối đúng”, mô hình dễ học đường tắt: miễn ra đáp án đúng, còn giải thích có thể “nghe hợp lý” nhưng sai.
            \item Nếu mỗi bước đều bị kiểm tra, mô hình bị phạt ngay khi “bịa bước”, nên có động lực tạo chuỗi bước nhất quán hơn.
        \end{itemize}
    \end{block}
    % Trong công trình "Let's Verify Step by Step", \cite{Lightman2023} đã so sánh hai kiểu mô hình phần thưởng, chứng minh được PRM giảm đáng kể được trường hợp “đúng đáp án nhưng sai lập luận”.
    \vspace{0.5cm}
    
    \begin{columns}[t]
        % Cột trái: ORM
        \begin{column}{0.48\textwidth}
            \centering
            \textbf{\textcolor{red}{Outcome Supervision (ORM)}}
            \vspace{0.2cm}
            
            \begin{tikzpicture}[scale=0.8, every node/.style={transform shape}]
                \node[draw, rounded corners] (start) {Input};
                \node[draw, rounded corners, right=2cm of start] (end) {Answer};
                \draw[->, thick, dashed] (start) -- node[above] {Blackbox} (end);
                \node[circle, draw, red, thick, right=0.2cm of end] {?};
            \end{tikzpicture}
        \end{column}
        
        % Vạch ngăn cách
        \vrule{}
        
        % Cột phải: PRM
        \begin{column}{0.48\textwidth}
            \centering
            \textbf{\textcolor{green!60!black}{Process Supervision (PRM)}}\footnote{Dựa trên bài báo: Let's Verify Step by Step (\cite{Lightman2023})}
            \vspace{0.2cm}
            
            \begin{tikzpicture}[scale=0.8, every node/.style={transform shape}]
                \node[draw, rounded corners] (s1) {Step 1};
                \node[circle, fill=green!20, right=0.1cm of s1, inner sep=1pt] {\tiny \checkmark};
                
                \node[draw, rounded corners, right=0.8cm of s1] (s2) {Step 2};
                \node[circle, fill=green!20, right=0.1cm of s2, inner sep=1pt] {\tiny \checkmark};
                
                \draw[->] (s1) -- (s2);
            \end{tikzpicture}
        \end{column}
    \end{columns}
    
    \vspace{0.5cm}
\end{frame}